% Root independent proof/code audit passed on 2026-11-15.
\begin{statement}{C.36}{Proved}{A nonzero arithmetic-trivial character at every prime at least thirteen}
The lower bound $d_{p-1}$ in C.34 is positive for every
prime $p\geq13$.  Hence the full gentle arithmetic-normal
quotient has a nonzero continuous character to
$\mathbb F_p$ for every such prime, and its abelianization
has the simultaneous product quotient of C.34 with
nontrivial factors at all these primes.
No absence assertion is made for smaller odd primes.
\end{statement}
\begin{proof}
We use the exact quantities $a_n,b_n,d_n$ of C.34.
The independent program
\texttt{pbc\_gentle\_uniform\_prime\_bound.py}
uses a Mobius sieve and the proved character formula
to compute all relevant integers.
It agrees with the previous implementation through
degree $300$ and verifies
\[
                 a_n>b_n\quad
                   (n=12,14,\ldots,998).                        \tag{C36.1}
\]
All $494$ exact rows are retained.  The following
analytic estimate proves every remaining degree,
rather than extrapolating this finite check.

Put $r_n=a_n2^{-n}$.  The trace formula implies, for
$n\geq2$,
\[
 \left|r_n-\frac1{3n(n+1)}\right|\leq3\,2^{-n/2}.                \tag{C36.2}
\]
Indeed the nonleading divisors in $L_n(1)$ contribute
at most $2^{n/2}$; the transposition trace has absolute
value at most $2^{n/2}$, and the three-cycle trace
has absolute value at most $2^{n/3}+1$.
Substitution in (C34.1), including the terms at $n+1$,
gives (C36.2).
For every $n\geq100$,
\[
 3\,2^{-n/2}\leq\frac1{100n^2}.
                                                               \tag{C36.3}
\]
The inequality holds at $n=100$, and
$n^2 2^{-n/2}$ decreases thereafter, since
$(101/100)^4<2$.
Combining these estimates gives the convenient rational
bounds
\[
             \frac8{25n^2}\leq r_n\leq\frac7{20n^2}
                         \qquad(n\geq100).                     \tag{C36.4}
\]
For the lower bound use
$100/303-1/100>8/25$; for the upper bound use
$1/3+1/100<7/20$.

By definition of the unordered bracket count,
\[
 \frac{b_n}{a_n}
    \leq\sum_{2\leq i\leq n/2}
                       r_i\,\frac{r_{n-i}}{r_n}.               \tag{C36.5}
\]
This slightly overcounts the middle term when $n$ is even,
which is harmless.
If $n\geq1000$ and $i\leq100$, (C36.4) gives
\[
 \frac{r_{n-i}}{r_n}
   \leq\frac{35}{32}\left(\frac{n}{n-i}\right)^2
   \leq\frac{875}{648}.
\]
The sum of these endpoint terms is at most
$(875/648)S$, where C.34 proves
$S=\sum r_i\leq7/16$.
For $100<i\leq n/2$, the same ratio is at most $35/8$,
and
\[
            \sum_{i>100}r_i
               \leq\frac7{20}\sum_{i>100}i^{-2}
               \leq\frac7{2000}.
\]
Consequently, for every $n\geq1000$,
\[
 \frac{b_n}{a_n}
 \leq \frac{875}{648}\frac7{16}
           +\frac{35}{8}\frac7{2000}
       =\frac{78547}{129600}<1.                                \tag{C36.6}
\]
Equations (C36.1) and (C36.6) prove $d_n>0$ at every
even $n\geq12$.  Set $n=p-1$ and apply the already
proved actual-character construction of C.34.
\end{proof}
